
%-------------------------------------------------------------------------------
\chapter{Calibration and Object/Event Definitions}
\label{sec:objects}
%-------------------------------------------------------------------------------


\todo{The calibration part is a summary from the W pT paper. This needs to be updated for the calibration of the mW}

\section{Electrons calibration and corrections}
\label{sec:elCorr}

Reference~\cite{lowMuMw:INT1} details the electron performance studies done using the special low-pileup datasets collected by the ATLAS detector at $\sqrt{s} = 5$ and 13~\TeV.
The electron reconstruction scale factors are obtained by extrapolation of the standard high-pileup SFs to the \lowmu regime, and apply to both the 13 and 5 \TeV{} datasets. The identification scale factors are measured in-situ separately for the 13 and 5 \TeV{} data.
Isolation and trigger efficiencies and SF are measured in-situ separately for the 13 and 5 \TeV{} data sets.
For the electron energy calibration, the global strategy follows the methodology applied in high-$\mu$ standard calibration and described in Ref.~\cite{EGAM-2018-01}.
Electron energy scale and resolution corrections are measured in-situ using \Zee events from the low pile-up dataset. 

\section{Muons calibration and correction}\label{sec:muCorr}

Reference ~\cite{lowMuMw:INT1} details the muon-related performance measurements obtained using the ATLAS low pile-up $\sqrt{s} = 5$
TeV and $13$ TeV data relevant for this measurement.
Muon reconstruction and \textit{track-to-vertex association} requirements efficiencies measured
insitu with \Zmm\ events are found to be compatible with the
corresponding high-$\mu$ measurements and are measured to
permille-level precision. The muon trigger and isolation efficiencies
are measured insitu in the \lowmu\ data, for each dataset, to typically better than $<1\%$ precision, limited
by the size of the \Zmm\ data samples. Finally, the muon calibration
derived from the high-$\mu$ data has been cross-checked using both
\Zmm\ and \Jmumu\ events in the \lowmu\ datasets and has been found to be adequate at the
current level of systematics. A dedicated correction for
charge-dependent momentum biases, so-called sagitta correction, has
been derived using a combination of different methods using the $13$ TeV \lowmu dataset.


\section{Hadronic Recoil calibration and correction}\label{sec:recoilCorr}

The detailed procedure to calibrate the hadronic recoil is described in Ref.~\cite{lowMuMw:INT1}. It is briefly summarised here. The calibration is obtained as a function of \setue\ and \ptv\ of the boson, since the hadronic recoil is mainly sensitive to these. The correction is obtained in \Zboson\ events and extrapolated to the \Wboson\ events.

The procedure consists of three steps :
\begin{itemize}
	\item First, the \setue\ distribution in the Monte-Carlo should be well modelled and match that of the data. More precisely, it is crucial to model correctly the correlation between \setue\ and \ptv, since we want to have a good description of the activity as a function of our measured physics observable. This is achieved thanks to a 2-dimensional reweighting, obtained in \Zboson\ events. In the simulated \Wboson\ events, an additional reweighting of \setue\ in bins of \ut\ is applied. A further 1-dimensional reweighting is obtained for each process (\Wminus, \Wplus\ and \Zboson) to recover the initial underlying \pttruth\ spectrum.

	\item Second, the direction of the recoil is corrected, taking the projections on $x$ and $y$ axes of the recoil in \Zboson\ events, and correcting for any data to MC differences.

	\item Finally, response and resolution corrections are, once again, obtained in-situ in \Zboson\ events, where the
    parallel (\upar) and perpendicular (\uperp) components can be extracted in the data as a function of \setue\ and \ptll, and compared
    to the Monte-Carlo to extract corrective coefficients. 
\end{itemize}

This \Zboson\ boson based calibration is applied to \Wboson\ events ; uncertainties due to this extrapolation are included. A summary on these uncertainties, as well as their impact on the unfolded spectra, are discussed in section~\ref{subsec:uncsummary}.

\section{Object and Event Definitions}
\label{subsec:definitions}

\todo{This section needs to be checked for any errors}


In this analysis, electrons or muons are considered as
leptons. Electrons are reconstructed from energy clusters in the
calorimeter associated to a track, while muons are reconstructed from a
track that crosses both the inner detector and the muon chambers. Both electrons and muons are required to
have $\pT>25~\GeV$. Additionally, electrons should have $|\eta|<2.47$
and candidates in the crack region $1.37<|\eta|<1.52$ are
excluded. Muons candidates are required to be within $|\eta|<2.4$. The
electron candidates must satisfy the Medium LH identification criteria
and are required to be isolated, satisfying the track isolation
criterion $\texttt{ptcone20}/\min(\pt, 50\,\GeV) <0.1$.  The muon
candidates must satisfy the Medium identification criteria and are
also required to be isolated, satisfying the
$\texttt{ptcone20}/\min(\pt, 50\,\GeV) <0.1$ track isolation criterion.

Leptons are required to originate from the primary vertex. The
longitudinal impact parameter of each lepton track, defined as the
distance between the track and the primary vertex along the beam line
multiplied by the sine of the track $\theta$ angle, is required to be
less than $0.5$ mm ( $\Delta z_{0}\times \sin(\theta)
<0.5$ mm). Furthermore, the significance of the transverse impact
parameter, defined by the transverse impact parameter ($d_{0}$ ) of a
lepton track with respect to the beam line, divided by its estimated
uncertainty ($\sigma(d_{0})$ ), is required to satisfy for electrons
$|d_0|/\sigma(d_{0}) < 5$ and for muons $|d_0|/\sigma(d_{0}) < 3$.

Dedicated lepton calibrations and efficiency corrections are applied to the data and the MC used in the analysis following the Reference~\cite{lowMuMw:INT1} for electrons and muons. A short summary of the calibration strategy and the corresponding resulting uncertainties for the electrons and muons is described in Sections~\ref{sec:elCorr} and ~\ref{sec:muCorr}.

In the context of W boson production a high precision computation of the neutrino kinematics is mandatory, and this is done with the help of the \textit{hadronic recoil}. In proton-proton collisions there is a non-zero transverse momentum for vector boson production that can be thought of to originate from initial state gluon and quark radiation:
\begin{equation}
\vec{p}_{T}(W/Z)=\vec{p}_{T}^{\mathrm{lepton1}}+\vec{p}_{T}^{\mathrm{lepton2}}=-\sum \vec{p}_{T}^{\mathrm{ISR quark,gluons}} = -\vec{\ut} \,.
\end{equation}
Here $p_{T}(W/Z)$ denotes the transverse momentum of the $W$ or $Z$ boson and $p_{T}^{\mathrm{lepton}}$ denote the transverse momenta decay leptons (charged lepton or neutrino). The  \textit{hadronic recoil}, $\sum \vec{p}_{T}^{\mathrm{ISR quark,gluons}}$,  is the quantity which accounts for the transverse momenta any partons recoiling from the colour-neutral boson, and is denoted as $\vec{\ut}$.
The  \textit{hadronic recoil} has been used successfully for the W mass measurement at 7 \TeV and in the \ptw analysis, where it was proven to be a powerful tool to determine the neutrino transverse momentum in a more precise way than the ``standard'' TST (track soft term) based algorithm~\cite{STDM-2014-18}. This analysis, as with the \ptw analysis, uses \emph{particle flow objects} (PFOs) as input constituents to the building of \ut, and therefore exploits the full available topo-clusters and track information. 

Other observables related to the neutrino which are used within this analysis, especially for the calibration, are:

\textbf{MET} is the proxy for neutrino $\nu$ transverse momentum defined as 
\begin{equation}\label{eq:METandNeutrino}
\vec{\met} :=\vec{\pt}^{\nu}=-(\vec{\ut}+\vec{\pt}^\ell)
\end{equation}
where $\vec{\met}$ is the missing transverse energy. The transverse momenta of additional muons with $\pT>5~\GeV$ and satisfying the medium identification is added at the definition of \MET{} while their calorimeter energy deposit is subtracted from the event.

\textbf{\set} which is a quantity that represents the total event activity. \set\ is computed as the scalar sum of the \pt of all
PFOs and receives contributions from pileup, noise, and
underlying event. The resolution of the \ut\ measurement will typically grow
$\propto \set$. It is also strongly dependent on the vector
boson dynamics and will grow with $\ptv=\ut$. 

\textbf{\setue} is defined as $\setue = \set - \ut$ and is used to disentangle the effects from pileup, noise and UE from the vector boson dynamics. \setue\ may be thought of the event activity corrected for the ``directed'' recoil activity and represents more
directly the activity from underlying event, pileup, and hard
emissions beyond the first~\footnote{E.g. a $V+1$ jet event has
  $\setue \sim 0$ on parton level for any jet \pt and thus a finite
  $\setue > 0$ is due to additional activity.}.




